\begin{tikzpicture}[tfm,
      every node/.append style={execute at begin node={\begin{hyphenrules}{nohyphenation}}, execute at end node={\end{hyphenrules}}}, % sin guiones
      obj/.style={arqbase, fill=arq-datos, draw=uoc-azul, text width=40mm, minimum height=19mm,
                  inner sep=4pt, font=\sffamily\scriptsize, align=left, text=tinta},
      objalg/.style={obj, fill=nn-conv, double=nn-conv, double distance=1.2pt},
      num/.style={circle, fill=uoc-azul, text=white, font=\sffamily\scriptsize\bfseries,
                  inner sep=0pt, minimum size=5mm},
    ]
    % ---------- Fondo del MVP
    \fill[uoc-gris] (-2.55, 2.15) rectangle (12.55,-4.65);
    \node[anchor=north west, font=\sffamily\scriptsize\bfseries, text=uoc-azul] at (-2.35, 2.05)
      {Producto mínimo viable (MVP): objetivos principales, en orden de ejecución};

    % ---------- Objetivos
    \node[obj]    (o1) at (0,0)   {\textbf{OP1}\\[1pt] Identificar necesidades y requisitos de accesibilidad universal};
    \node[obj]    (o2) at (5,0)   {\textbf{OP2}\\[1pt] Identificar, analizar y preparar los datos};
    \node[obj]    (o3) at (10,0)  {\textbf{OP3}\\[1pt] Identificar y definir las métricas e indicadores de evaluación de la accesibilidad};
    \node[objalg] (o4) at (0,-3.3) {\textbf{OP4}\\[1pt] Entrenar modelos de aprendizaje automático de clasificación de accesibilidad};
    \node[obj]    (o5) at (5,-3.3) {\textbf{OP5}\\[1pt] Validación de los modelos};
    \node[obj]    (o6) at (10,-3.3){\textbf{OP6}\\[1pt] Interpretación de los resultados};

    % ---------- Números de orden
    \foreach \i in {1,...,6}
      \node[num] at (o\i.north west) {\i};

    % ---------- Flechas en orden
    \draw[flecha] (o1) -- (o2);
    \draw[flecha] (o2) -- (o3);
    \draw[flecha] (o3.south) -- ++(0,-0.62) -| (o4.north);
    \draw[flecha] (o4) -- (o5);
    \draw[flecha] (o5) -- (o6);
  \end{tikzpicture}

  \caption{Objetivos principales del proyecto (MVP).}
  \label{fig:objetivos_mvp}
  \fuente{elaboración propia, con ayuda de Claude.}
